\chapter{\nt{SERO}ニューロンは何を計算するか}
\label{sec:serocomputer}

\section{最初のブロードキャストチャネル}

これまで本書のニューロンはすべてアドレスバスで話していた。\nt{GLUT}と\nt{GABA}は決まった受け手にスパイクを送り、そのようなネットワークが何を計算し、どんな言語で話すかを私たちは調べてきた。ここからは\ref{sec:buses}節の第二のバス、ブロードキャストバスに移る。その最初のチャネルが\nt{SERO}である。これまでと同じく、生きた生体で実際に起こることだけを使う。

この章では、以後すべてのブロードキャストチャネルが使う三つの部品が登場する。一定の背景入力があれば止められるまで自分で動き続けるニューロン、実は組織の体積である配線、そして個々のスパイクの代わりに使うゆっくりした濃度である。第\ref{sec:dofa}--\ref{sec:acet}章の\nt{DOPA}、\nt{NORA}、\nt{ACET}も同じ部品から組み立てられる。

エンジニアの目で見ると、\nt{SERO}ニューロンには重要な違いが四つある。

\emph{第一に、符号は受け手が選ぶ。}セロトニンには両方の符号の受容体があり、規則~\eqref{eq:dalesign}はここでは成り立たない。符号を決めるのは送り手ではなく、受け手の受容体である（命題~\ref{prop:receptor}）~\cite{BarnesSharp1999}。

\emph{第二に、一定の背景があれば\nt{SERO}ニューロンは自分で動く。}覚醒時には毎秒数回、一定のリズムでスパイクを出し、スパイクごとの後押しを必要としない。つまりペースメーカーである。ただし一定の背景入力は欠かせない。この入力がない脳スライスでは、こうした細胞の大半は沈黙している。ノルアドレナリンを$\alpha_1$受容体経由で与えると一定のリズムが戻り、その受容体を遮断すると消える~\cite{VandermaelenAghajanian1983}。生体ではこの背景は\nt{NORA}から来る。回路として見れば電源を要する発振器であり、電源がある限り、抑制で止められるまで自分で動く。

\emph{第三に、遅い。}セロトニン受容体はほぼすべて代謝型である。自分でチャネルを開くのではなく、細胞内で一連の反応を起動する。応答は遅く、主要な受容体を通る抑制電流はおよそ1秒のうちに立ち上がって減衰する~\cite{CourtneyFord2016}。これは\nt{GLUT}の速いシナプスの応答より数十倍長い。

\emph{第四に、神経伝達物質を配線ではなく体積に放出する。}\nt{SERO}ニューロンの出力が届く領域では、セロトニンは狭い間隙ではなく細胞外空間に出て、周囲に広がることが多い~\cite{Bunin1998}。受け手が感じるのは濃度であり、分子がどの送り手から来たかは知りようがない。

このような時間スケールでは個々のスパイクはもはや意味をもたないので、ネットワークは発火率$r_j$と濃度で記述する。ニューロン$j$が伝達物質を放出する体積内のセロトニン濃度$c$は、放出によって増え、伝達物質が回収されることで減る。
\begin{empheq}[box=\keybox]{equation}
\tau_c\,\frac{\dd c}{\dd t} \;=\; -c \;+\; \kappa\sum_j r_j ,
\qquad
r_i \;=\; f_i\bigl(b_i + s_i\,g\,c\bigr), \quad s_i=\pm1 .
\label{eq:seroneuron}
\end{empheq}
これはスパイク列のもう一つの読み方であり、命題~\ref{prop:reader}を補う。体積には個々のスパイクもその順序も見えず、時間$\tau_c$で平均した発火率だけが見える。第一式は膜と同じRC回路である。$\tau_c$は体積内での伝達物質の寿命で、直接には測定されておらず、ここでは1秒程度とおく。$\kappa$は1スパイクあたりの伝達物質の量である。第二式は受け手を記述する。$b_i$は受け手自身の入力で、ペースメーカーでは正である（\nt{NORA}からの一定の背景入力もここに含まれる）。$g$は受容体の感度、符号$s_i$は受け手が選び、$f_i$は非減少の伝達特性である。

\section{1ニューロンでNOT}

抑制性受容体（$s=-1$）をもつペースメーカーを考える。周囲にセロトニンがなければ、自身の入力$b$で動く。セロトニンが現れると入力は$gc$だけ減り、$gc>b$ならニューロンは沈黙する。これがNOTである。入力がないときに出力がある。使ったのはニューロン1個である（図~\ref{fig:serogates}）。命題~\ref{prop:monotone}はここには当てはまらない。あれは正の重みを前提としていた。

同じペースメーカーに二つの異なる体積からセロトニンが作用すれば、少なくとも一方にセロトニンがあるとき沈黙する。これはNORであり、NORからは任意の論理関数を組み立てられる。\nt{GLUT}ネットワークの2線式符号はここでは要らない。信号がないことは、止められるまで動き続けるニューロンが計算する。

\begin{figure}[t]
\centering
\begin{tikzpicture}[>=Stealth,
  nrn/.style={circle,draw,very thick,fill=blue!6,minimum size=0.95cm,inner sep=0pt},
  pace/.style={nrn,double,double distance=1.2pt},
  inh/.style={-{Bar[width=7pt]},thick,red!70!black}]
\node[pace] (N1) at (0,0) {};
\draw[inh] (-1.6,0) node[left,black] {$c$} -- (N1);
\draw[->,very thick,red!70!black] (N1) -- (1.3,0) node[right,black] {$\bar c$};
\node[below=0.55cm] at (0,0) {\small NOT};
\node[pace] (N2) at (5,0) {};
\draw[inh] (3.4,0.45) node[left,black] {$c_1$} -- (N2);
\draw[inh] (3.4,-0.45) node[left,black] {$c_2$} -- (N2);
\draw[->,very thick,red!70!black] (N2) -- (6.3,0) node[right,black] {$\overline{c_1\vee c_2}$};
\node[below=0.55cm] at (5,0) {\small NOR};
\end{tikzpicture}
\caption{\nt{SERO}ニューロンによる論理。二重丸はペースメーカーで、一定の背景入力があれば抑制されるまで自分で動く。横棒付きの線は抑制性受容体（$s=-1$）。$c$、$c_1$、$c_2$は受け手の周囲の体積におけるセロトニン濃度。左がNOT、右がNOR。}
\label{fig:serogates}
\end{figure}

\begin{keyblock}
\begin{proposition}[\nt{SERO}論理の完全性]
\label{prop:serocomplete}
ネットワークに抑制性受容体をもつペースメーカーがあり、異なる体積への放出が混ざらないなら、ネットワーク~\eqref{eq:seroneuron}は任意の論理関数を計算でき、各ビットは1本の配線で伝わる。
\end{proposition}
\end{keyblock}

論理の面では\nt{SERO}ネットワークは\nt{GLUT}より強い。しかし「異なる体積への放出が混ざらない」という条件は脳では満たされない（\ref{sec:volume}節）。

\section{負帰還}

セロトニンニューロンは、自分自身の上に抑制性受容体をもつ~\cite{Sprouse1987}。放出したセロトニンが自分に戻ってきて、自分にブレーキをかける。エンジニアにはおなじみの回路、負帰還増幅器である（図~\ref{fig:serofeedback}）。

何が測定済みで、何がモデルかを区別しておく。\nt{SERO}ニューロンが集まる縫線核そのものの中では、セロトニンは共通の体積ではなく、シナプスを通じて一対一で近隣を抑制する。トランスポーターがすばやく回収し、間隙の外に溜まらせないからである。1回の放出は発火に短い休止を生むだけで、平均発火率は変わらない~\cite{CourtneyFord2016}。したがって以下の、共通の体積を通じて閉じたループはモデルである。このような帰還が平均発火率のレベルで働いたら何をするかを計算する。

\begin{figure}[t]
\centering
\begin{tikzpicture}[>=Stealth,
  nrn/.style={circle,draw,very thick,fill=blue!6,minimum size=0.95cm,inner sep=0pt},
  pace/.style={nrn,double,double distance=1.2pt},
  blk/.style={draw,very thick,rounded corners=2pt,fill=orange!10,minimum height=0.9cm,minimum width=2.2cm,align=center,font=\small},
  inh/.style={-{Bar[width=7pt]},thick,red!70!black}]
\node[pace] (N) at (0,0) {};
\node[blk] (V) at (3.6,0) {体積\\$\tau_c$};
\draw[->,very thick] (N) -- node[above] {\small $r$} (V);
\draw[->,very thick] (V) -- (6.0,0) node[right] {\small $c$：すべての受け手へ};
\draw[inh] (V.south) -- ++(0,-0.9) -| node[pos=0.25,below] {\small $-g\,c$} (N.south);
\draw[->,thick,blue!60!black] (-1.7,0) node[left,black] {\small $b$} -- (N);
\end{tikzpicture}
\caption{自身の伝達物質による帰還をもつ\nt{SERO}ニューロン。濃度$c$はすべての受け手に届き、ニューロン自身も抑制する。モデルではこれはレベル調節器である。脳でループがこの役割を果たすというのは仮説である。}
\label{fig:serofeedback}
\end{figure}

このループが何をするかを計算しよう。伝達特性を線形とし、$u>0$で$f(u)=au$とする。平衡では$c=\kappa r$かつ$r=a(b-gc)$である。一方を他方に代入すると
\begin{empheq}[box=\keybox]{equation}
r \;=\; \frac{a\,b}{1+G}, \qquad G \;=\; a\,g\,\kappa .
\label{eq:feedback}
\end{empheq}
$G$はループゲインである。これはどんな負帰還増幅器にも共通の式であり、同じ二つの利点をもたらす。

\emph{ばらつきへの強さ。}ニューロンの興奮性$a$が割合$\varepsilon$だけ変わったとする。帰還がなければ発火率も同じ割合で変わる。帰還があり$G\gg1$なら、発火率$r\approx b/(g\kappa)$は$a$にほとんど依存せず、その変化は$1+G$分の1に縮む。$G=9$ならばらつきは10分の1に抑えられる。

\emph{速さ。}$r=a(b-gc)$を式~\eqref{eq:seroneuron}の第一式に代入すると$\tau_c\,\dd c/\dd t=-(1+G)\,c+\kappa ab$となる。濃度は時定数$\tau_c/(1+G)$で自分のレベルに近づく。帰還がない場合の$1+G$倍速い。

これが\nt{SERO}系が担いうる最初の計算である。サーモスタットが温度を保つように、脳全体のセロトニンのレベルを設定値に保つ。これは仮説である。実験では、自己抑制は今のところ平均発火率を変えない短い休止としてしか見えていない。

\section{スパイクからレベルへ}

第二の計算は式~\eqref{eq:seroneuron}の第一式に隠れている。ニューロンは個々のスパイクを出すが、受け手が感じるのは、それらを時間$\tau_c$で平滑化した濃度である。源の発火率が変わると、濃度は遅れてそれに追従する。
\begin{equation}
c(t) \;=\; \frac{\kappa}{\tau_c}\int_{-\infty}^{t} e^{-(t-t')/\tau_c}\,\sum_j r_j(t')\,\dd t' .
\label{eq:lowpass}
\end{equation}
これは直近の数$\tau_c$にわたる源の活動の移動平均、すなわちローパスフィルタである（図~\ref{fig:lowpass}）。最も簡単なD/A変換器も、パルスをRC回路に通してその頻度をレベルに変える。\nt{SERO}系は同じことを数十万個のニューロンで一度に行う。

この量は同時に記憶でもある。源が沈黙した後も、濃度はさらに$\tau_c$程度の時間保たれる。これはビットではなく、なめらかに消えていく痕跡である。

\begin{figure}[t]
\centering
\begin{tikzpicture}[>=Stealth,x=1.45cm,y=1cm]
\draw[gray!70!black] (0,3.2) -- (8.1,3.2);
\node[left,font=\small] at (-0.05,3.4) {スパイク};
\node[above,font=\scriptsize,gray!40!black] at (1,3.65) {毎秒$2$回};
\node[above,font=\scriptsize,gray!40!black] at (3.5,3.65) {毎秒$8$回};
\node[above,font=\scriptsize,gray!40!black] at (6.5,3.65) {毎秒$2$回};
\draw[->,gray!70!black] (0,0) -- (8.3,0) node[right,black] {\small $t$};
\draw[->,gray!70!black] (0,0) -- (0,2.75) node[left,black] {\small $c$};
\foreach \s in {0.136,0.529,0.760,1.223,1.714,1.748,1.755,2.663,2.701,2.734,3.414,3.493,3.719,3.800,3.928,3.948,4.074,4.327,4.420,4.589,4.728,4.736,4.914,5.025,5.205,5.220,6.224,6.544,7.178}{\draw[thick,blue!60!black] (\s,3.2) -- (\s,3.6);}
\foreach \a/\b/\v in {0.000/0.136/0.000,0.136/0.529/1.000,0.529/0.760/1.675,0.760/1.223/2.330,1.223/1.714/2.466,1.714/1.748/2.509,1.748/1.755/3.425,1.755/2.663/4.402,2.663/2.701/2.775,2.701/2.734/3.672,2.734/3.414/4.553,3.414/3.493/3.306,3.493/3.719/4.055,3.719/3.800/4.235,3.800/3.928/4.905,3.928/3.948/5.316,3.948/4.074/6.211,4.074/4.327/6.476,4.327/4.420/6.028,4.420/4.589/6.493,4.589/4.728/6.483,4.728/4.736/6.642,4.736/4.914/7.589,4.914/5.025/7.351,5.025/5.205/7.579,5.205/5.220/7.331,5.220/6.224/8.221,6.224/6.544/4.012,6.544/7.178/3.914,7.178/8.000/3.076}{\draw[very thick,orange!85!black] plot[domain=\a:\b,samples=10] (\x,{0.25*\v*exp(-(\x-\a))});}
\foreach \s/\p/\q in {0.136/0.000/1.000,0.529/0.675/1.675,0.760/1.330/2.330,1.223/1.466/2.466,1.714/1.509/2.509,1.748/2.425/3.425,1.755/3.402/4.402,2.663/1.775/2.775,2.701/2.672/3.672,2.734/3.553/4.553,3.414/2.306/3.306,3.493/3.055/4.055,3.719/3.235/4.235,3.800/3.905/4.905,3.928/4.316/5.316,3.948/5.211/6.211,4.074/5.476/6.476,4.327/5.028/6.028,4.420/5.493/6.493,4.589/5.483/6.483,4.728/5.642/6.642,4.736/6.589/7.589,4.914/6.351/7.351,5.025/6.579/7.579,5.205/6.331/7.331,5.220/7.221/8.221,6.224/3.012/4.012,6.544/2.914/3.914,7.178/2.076/3.076}{\draw[very thick,orange!85!black] (\s,0.25*\p) -- (\s,0.25*\q);}
\draw[thick,densely dashed,black] plot[domain=0:2,samples=30] (\x,{0.25*2*(1-exp(-\x))})
  plot[domain=2:5,samples=40] (\x,{0.25*(8-6.271*exp(-(\x-2)))})
  plot[domain=5:8,samples=40] (\x,{0.25*(2+5.688*exp(-(\x-5)))});
\foreach \x in {0,2,5,8}{\draw[gray!60!black] (\x,-0.05) -- (\x,0.05) node[below=3pt,font=\scriptsize] {\x\,s};}
\end{tikzpicture}
\caption{スパイクからレベルへ。上は\nt{SERO}ニューロンのスパイク。2秒間はまばら、3秒間は密、その後ふたたびまばらになる。下は$\tau_c=1\,\text{s}$のときのセロトニン濃度~\eqref{eq:lowpass}。破線はスパイクが完全に等間隔だった場合の同じ量。}
\label{fig:lowpass}
\end{figure}

\section{配線とは体積である}
\label{sec:volume}

命題~\ref{prop:serocomplete}の条件に戻ろう。近くで別々の送り手が放出したセロトニンは、一つの濃度に足し合わされる。

\emph{ここでの配線は線維ではなく、組織の体積である。}二つの信号が別々のままでいられるのは、伝達物質が広がりきる距離よりも離れた領域に放出される場合だけである。拡散係数$D$の分子は、時間$\tau$のうちに
\begin{empheq}[box=\keybox]{equation}
\ell \;\sim\; \sqrt{D\,\tau}\, .
\label{eq:diffusionlength}
\end{empheq}
だけ移動する。細胞間隙の屈曲は小分子の拡散をおよそ$2.6$倍遅くする~\cite{Sykova2008}。セロトニン自体の組織内拡散係数は測定されていない。分子の大きさから、数$\times10^{-6}$~cm$^2$/sと見積もる。信号の寿命が$\tau\sim1$~s（これも推定）なら、$\ell\sim\sqrt{3\cdot10^{-6}}$~cm $\approx20$~\textmu mとなる。このネットワークでのアドレスは\emph{場所}である。受け手は、自分がどこにいるかによってしか、信号が誰から来たかを知ることができない。

実際の\nt{SERO}ニューロンは、こうした個別のアドレスがほとんど残らないように作られている。それぞれの出力は脳の広大な範囲にばらまかれ、1個の細胞の枝が互いに遠く離れた多くの領域に届く~\cite{GagnonParent2014}。1細胞あたりの放出部位は推定でおよそ$10^5$（表~\ref{tab:types}）で、直接数えられたわけではない。異なる\nt{SERO}ニューロンの「配線」は重なり合い、混ざる。それでも完全に1本の配線に融合するわけではない。\nt{SERO}ニューロンはいくつかのサブシステムに分かれ、それぞれ異なる領域に出力を送り、同じ出来事に異なる応答をする~\cite{Ren2018}。しかしそうしたサブシステムの数は数個であって、数十万ではない。

\begin{keyblock}
\begin{proposition}[独立な配線の数]
\label{prop:volumewires}
容積伝達のネットワークでは、独立な信号の数を制限するのはニューロンの数ではなく、ニューロンが伝達物質を放出する、大きさ$\ell\sim\sqrt{D\tau}$の重ならない領域の数である。各ニューロンの出力が大きな体積にばらまかれていれば、ネットワーク全体は少数の共通変数しか伝えない。
\end{proposition}
\end{keyblock}

したがって、命題~\ref{prop:serocomplete}の論理回路は、脳では組み立てる材料がない。一方、調節器~\eqref{eq:feedback}とフィルタ~\eqref{eq:lowpass}には個別の配線は要らず、共通の量が一つあれば足りる。フィルタは投射先領域での体積への放出の測定から直接導かれる~\cite{Bunin1998}。レベル調節器は仮説である。

\section{計画の時間幅}
\label{sec:horizon}

共通でゆっくりした量が一つあると、何の役に立つだろうか。よい候補は、ネットワーク全体で共通であるべきで、数秒から数分より速くは変わらないパラメータである。そのようなパラメータは第\ref{sec:neurontypes}章ですでに登場した。誤差の式~\eqref{eq:rpe}の係数$\gamma$で、将来の報酬が即時の報酬よりどれだけ安いかを表す。連続時間ではその役割を\emph{時間幅}$\tau_\gamma$が担う。時間$D$待たなければならない報酬$R$は、いま$R\,e^{-D/\tau_\gamma}$の価値をもつ。

時間幅は待つ価値があるかどうかを決める。小さな報酬$r_0$をすぐ取るか、大きな報酬$R$を待つかを選べるとする。待つほうが得なのは
\begin{empheq}[box=\keybox]{equation}
D \;<\; \tau_\gamma\,\ln\frac{R}{r_0} .
\label{eq:patience}
\end{empheq}
の間である。$\tau_\gamma=2\,\text{s}$で遅延報酬が3倍なら、$2.2\,\text{s}$まで待つ価値がある。時間幅が2倍なら$4.4\,\text{s}$までである。忍耐は時間幅に比例する。

式~\eqref{eq:patience}は指数割引に基づいている。秒単位の遅延で測定された割引は、双曲線$R/(1+D/\tau_\gamma)$に近い。サルでは、選択における遅延報酬の価値も、その予告信号に対する\nt{DOPA}ニューロンの応答も、このように下がる~\cite{KobayashiSchultz2008}。双曲線では条件~\eqref{eq:patience}は$D<\tau_\gamma\,(R/r_0-1)$に置き換わる。報酬の比への依存は対数ではなく線形になり、同じ数値で待機の限界はほぼ2倍に伸びる。遅延がランダムなら同じ指数関数から双曲線が得られ（問題~\ref{ex:05:7}）、忍耐が時間スケールに比例することは変わらない。

仮説によれば、この時間幅を設定するのが\nt{SERO}である~\cite{Doya2002}。その役割に必要なものはそろっている。ネットワーク全体で一つの共通レベルがあり、それが数秒で変わる。直接の実験もある。セロトニンニューロンを人工的に活性化すると、動物は遅延報酬をあきらめるまでより長く待ち~\cite{Miyazaki2014}、報酬がまれになった場所でもより長く報酬を得ようとし続ける~\cite{Lottem2018}。セロトニン濃度がネットワーク内部でどのように$\tau_\gamma$に変わるのかは分かっていない。次章では、\nt{DOPA}が送り出す誤差に時間幅が入る。

\section{まとめ}

\begin{table}[t]
\centering
\small
\begin{tabular}{>{\raggedright}p{3.6cm}>{\raggedright}p{4.3cm}>{\raggedright\arraybackslash}p{4.3cm}}
\toprule
& \nt{GLUT} & \nt{SERO} \\
\midrule
反応時間 & ミリ秒 & 1秒未満から1秒 \\
作用の符号 & $+$のみ & $+$と$-$、受け手が選ぶ \\
入力がないとき & 沈黙 & 一定の背景入力があれば自分で動く \\
アドレス指定 & 一対一 & 体積、アドレスは場所 \\
NOT & 「オフ」検出器経由のみ & ニューロン1個 \\
記憶 & 自己維持する活動、疲労で消える & 濃度、時間$\tau_c$で消える \\
主な計算 & 同時性、遅延、論理、系列 & 平滑化、レベル調節と時間幅$\tau_\gamma$（仮説） \\
\bottomrule
\end{tabular}
\caption{\nt{GLUT}ニューロンだけのネットワークと\nt{SERO}ニューロンだけのネットワークが計算するもの。}
\label{tab:glutvssero}
\end{table}

論理素子としては（表~\ref{tab:glutvssero}）\nt{SERO}ニューロンは\nt{GLUT}より豊かである。しかし通信系としては一斉同報であり、遅く、アドレスもほとんどない。だからプロセッサになろうとはせず、第\ref{sec:neurontypes}章で述べた少数の共通量を平滑化して配る。仮説によれば、計画の時間幅もその一つである。ペースメーカー、体積、ゆっくりした濃度には、この後さらに3回、\nt{DOPA}、\nt{NORA}、\nt{ACET}で出会う。

\section{問題}

\begin{exercise}[\lvlA]
\label{ex:05:1}
\emph{配線とは体積である。}セロトニンについて$D=3\cdot10^{-6}$~cm$^2$/sとする（\ref{sec:volume}節）。\eqref{eq:diffusionlength}から、寿命$1\,\text{s}$の信号の$\ell$と、伝達物質が$5$~cm広がるのにかかる時間を求めよ。では\nt{SERO}のブロードキャスト性を支えているのは、拡散か、それとも放出部位$\sim10^5$の配線の分岐か。
\end{exercise}

\begin{exercise}[\lvlA]
\label{ex:05:2}
\emph{レベル調節器。}\eqref{eq:seroneuron}で$f(u)=au$のとき、相対感度$S=(a/r)\,\partial r/\partial a$が$G\gg1$に限らず厳密に$1/(1+G)$に等しいことを示せ。$G=9$で興奮性$a$が$20\,\%$増えた。線形化せずに求めると、$r$は何パーセント変わるか。
\end{exercise}

\begin{exercise}[\lvlB]
\label{ex:05:3}
\emph{ループ内の遅い受容体。}\nt{SERO}受容体は遅れて応答する。$r(t)=a\bigl(b-gc(t-T)\bigr)$、体積は\eqref{eq:seroneuron}に従う。$G>1$のとき、ループが安定なのは$T<T^*=\tau_c\arccos(-1/G)/\sqrt{G^2-1}$の場合に限ることを示せ。$\tau_c=1\,\text{s}$、$G=2$と$G=9$について$T^*$を求めよ。遅延が数百ミリ秒のとき、どれだけのばらつき抑制が可能か。
\end{exercise}

\begin{exercise}[\lvlB]
\label{ex:05:4}
\emph{レベルのショットノイズ。}各スパイクは\eqref{eq:lowpass}に従い、$c$に$\tau_c=1\,\text{s}$の指数関数を加える。$3\cdot10^5$個の\nt{SERO}ニューロンすべてが毎秒$2$回のポアソンスパイクを一つの体積に放出するとき、$c$の変動係数を求めよ。毎秒$4$スパイクのニューロン1個で$\mathrm{CV}=10\,\%$を得るには、$\tau_c$はいくらであればよいか。
\end{exercise}

\begin{exercise}[\lvlB]
\label{ex:05:5}
\emph{シミュレーション：帰還対ノイズ。}発火率$r_i=a_i[b-gc]_+$（$a_i$は$[2.8,\,5.2]$で一様）の$N=50$個のポアソン型ペースメーカーと、$\kappa=1/N$、$\tau_c=1\,\text{s}$の体積~\eqref{eq:seroneuron}をシミュレートせよ。帰還（$b=1$、$g=2.25$）は、$b=0.1$、$g=0$の場合に比べてレベル$c$のCVを何分の1にするか。ニューロンの発火率はそろうか。
\end{exercise}

\begin{exercise}[\lvlB]
\label{ex:05:6}
\emph{設計：\nt{SERO}論理によるXOR。}ゲートは、$b-g\sum_kc_k>0$なら$r_0$を、そうでなければ$0$を自分の体積$\tau_c\,\dd c/\dd t=-c+\kappa r$に出すペースメーカーで、NORを計算する。このゲート5個でXORを組み、$\tau_c=1\,\text{s}$、$G'=g\kappa r_0/b=2$での整定時間を求めよ。
\end{exercise}

\begin{exercise}[\lvlB]
\label{ex:05:7}
\emph{遅延が未知のときの忍耐。}(a)~動物は、待ち時間が$6\,\text{s}$以内なら3倍の報酬を待つことに応じる。これはどんな時間幅$\tau_\gamma$を意味するか。(b)~遅延$D$が平均$\bar D$の指数分布に従うとする。待つほうが得なのは$\bar D<\tau_\gamma(R/r_0-1)$のときであることを示し、$\tau_\gamma=2\,\text{s}$、$R=3r_0$で固定遅延の場合と比べよ。
\end{exercise}

\begin{exercise}[\lvlC]
\label{ex:05:8}
\emph{濃度はどのように時間幅を設定するか。}\nt{DOPA}ニューロンは$V$を、直接には重み$k_1$で、$\tau_d=0.1\,\text{s}$で平滑化する\nt{GABA}仲介ニューロン経由では重み$-k_2$で受け取る。ゆっくりした$V$では和は$(k_1-k_2)V+k_2\tau_d\dot V$となる。$\tau_\gamma=2\,\text{s}$の\eqref{eq:ctd}に合うように$k_1$、$k_2$を選べ。\nt{SERO}は$k_1$を$m$倍する。時間幅を2倍にするには$m$をどれだけ変えればよいか。
\end{exercise}
